\subsection{Uniqueness checking}
\label{sec:uniqueness_checking}
\begin{sloppypar}
% Tiêu đề \paragraph nằm trên một dòng riêng; sloppypar là một nhóm nên cấu
% hình này chỉ áp dụng trong mục này.
\titleformat{\paragraph}[hang]{\normalfont\bfseries}{}{0pt}{}
\titlespacing*{\paragraph}{0pt}{8pt}{2pt}

Một bản đồ Griductive hợp lệ phải có đúng một lời giải, bởi nếu tồn tại hai cấu
hình tội trạng cùng thỏa mọi lời khai thì không có cách nào phân biệt chúng bằng
suy luận, và người chơi buộc phải đoán. Trong hệ thống, tính duy nhất của nghiệm
được kiểm tra bởi hàm \texttt{\_check\_unique\_solution} trong
\texttt{BackEnd/GameEngine/validation.py}. Hàm này là bước thứ hai của
\texttt{validate\_puzzle}, chạy sau bước kiểm tra cấu trúc (và kiểm tra mỗi ô có
đúng một lời khai nếu bật \texttt{require\_clue\_per\_cell}), đồng thời chạy
trước bước kiểm tra khả năng giải không cần đoán. Cũng giống engine, hàm chỉ
làm việc với bộ giải thông qua giao diện \texttt{ConstraintChecker}, nên có thể
dùng \texttt{DPLLChecker} hoặc \texttt{BruteForceChecker} tùy ý.

\paragraph{Tri thức dùng để kiểm tra}
Khác với ván chơi, nơi bộ giải chỉ thấy $\mathit{KB}_t$, bước kiểm tra này được
phép dùng tri thức đầy đủ của người thiết kế, tức toàn bộ tập lời khai
$\mathcal{K}$ kể cả những lời khai chưa lật (Mục~\ref{sec:knowledge_base}). Hàm
dựng tri thức đó bằng một hàm cục bộ \texttt{all\_clues\_state(proved)}, trả về
một \texttt{PublicState} có \texttt{revealed\_clues} là \emph{toàn bộ}
\texttt{puzzle.clues} và \texttt{proved\_statuses} là phép gán được truyền vào.
Như vậy, với một tập phán quyết $P$ cho trước, trạng thái này tương ứng với công
thức
\begin{equation}
  \mathit{KB}_{\mathrm{full}}(P) \;=\;
  \bigwedge_{\gamma \in \mathcal{K}} \phi_\gamma \;\land\;
  \bigwedge_{\langle i, s\rangle \in P} \phi_{\mathrm{Fact}(i,s)} .
  \label{eq:kb_full}
\end{equation}
Đây cũng là lý do \texttt{validate\_puzzle} là nơi duy nhất trong hệ thống tự
dựng \texttt{PublicState} mà không đi qua \texttt{GameEngine}
(Mục~\ref{sec:information_visibility}).

\paragraph{Bước 1: đáp án khai báo thỏa toàn bộ lời khai}
Hàm trước hết gọi \texttt{has\_solution} trên
$\mathit{KB}_{\mathrm{full}}(\sigma^\star)$, trong đó toàn bộ đáp án khai báo
$\sigma^\star$ được đưa vào dưới dạng các ô đã chứng minh và không kèm giả thiết
nào. Vì mọi biến ô đều đã bị cố định, truy vấn này thực chất chỉ kiểm tra
$\sigma^\star$ có thỏa mọi lời khai hay không: với \texttt{BruteForceChecker}
không còn ô tự do nào để duyệt, còn với \texttt{DPLLChecker} các mệnh đề đơn vị
của phán quyết được lan truyền ngay từ đầu. Nếu truy vấn không thỏa được, hàm
dừng với lỗi \texttt{Declared solution does not satisfy all clues}, và ca kiểm
thử \texttt{test\_declared\_solution\_violates\_clues} kiểm tra đúng trường hợp
này bằng cách đổi chỗ tội phạm trong đáp án.

\paragraph{Bước 2: không tồn tại nghiệm thứ hai}
Sau khi biết $\sigma^\star$ là một nghiệm, hàm dựng trạng thái gốc
$\mathit{KB}_{\mathrm{full}}(\mathcal{P}_0)$, trong đó các phán quyết ban đầu
\texttt{initial\_verdicts} được cố định cùng toàn bộ lời khai. Với mỗi ô $i$
không thuộc $\mathcal{P}_0$, hàm thử đảo trạng thái của ô đó so với đáp án bằng
truy vấn
\begin{equation}
  \mathrm{SAT}\bigl(\mathit{KB}_{\mathrm{full}}(\mathcal{P}_0) \land
  (x_i \ne \sigma^\star(i))\bigr),
  \label{eq:flip_query}
\end{equation}
được cài đặt là \texttt{has\_solution(base, Verdict(cell\_id, not
is\_criminal))}. Mỗi ô làm truy vấn này thỏa được sẽ được đưa vào danh sách
\texttt{ambiguous}. Hàm không dừng ở ô đầu tiên mà duyệt hết mọi ô rồi mới báo
lỗi \texttt{Solution is not unique: cells [...] can differ} kèm danh sách đã sắp
xếp, nhờ đó người thiết kế biết ngay tất cả những ô còn mơ hồ.
Thuật toán~\ref{alg:unique_solution} tóm tắt cả hai bước.

\begin{algorithm}[htbp]
  \caption{\texttt{\_check\_unique\_solution}}
  \label{alg:unique_solution}
  \begin{algorithmic}[1]
    \Require bản đồ với tập lời khai $\mathcal{K}$, đáp án $\sigma^\star$,
      phán quyết ban đầu $\mathcal{P}_0$; bộ kiểm tra \textsc{HasSolution}
    \Ensure không trả gì nếu nghiệm duy nhất, ngược lại báo lỗi
    \If{\textbf{not} $\textsc{HasSolution}(\mathit{KB}_{\mathrm{full}}(\sigma^\star))$}
      \State \textbf{raise} ``đáp án không thỏa toàn bộ lời khai''
    \EndIf
    \State $\mathit{ambiguous} \gets [\,]$
    \For{mỗi ô $i \notin \mathcal{P}_0$}
      \If{$\textsc{HasSolution}(\mathit{KB}_{\mathrm{full}}(\mathcal{P}_0),\ \langle i, \neg\sigma^\star(i)\rangle)$}
        \State thêm $i$ vào $\mathit{ambiguous}$
      \EndIf
    \EndFor
    \If{$\mathit{ambiguous} \ne [\,]$}
      \State \textbf{raise} ``nghiệm không duy nhất'' kèm $\mathit{ambiguous}$
    \EndIf
  \end{algorithmic}
\end{algorithm}

\paragraph{Vì sao đảo từng ô là đủ}
Hai bước trên tương đương với khẳng định $\sigma^\star$ là mô hình duy nhất của
$\mathit{KB}_{\mathrm{full}}(\mathcal{P}_0)$ khi xét trên các biến ô. Thật vậy,
nếu tồn tại một mô hình $\sigma' \ne \sigma^\star$ thì $\sigma'$ phải khác
$\sigma^\star$ tại ít nhất một ô $i$. Ô này không thể thuộc $\mathcal{P}_0$, vì
mọi mô hình đều phải khớp với các phán quyết ban đầu, còn các phán quyết ban đầu
đã được bước kiểm tra cấu trúc xác nhận là trùng với đáp án. Khi đó $\sigma'$
chính là nghiệm của truy vấn \eqref{eq:flip_query} tại ô $i$, nên ô $i$ bị đánh
dấu mơ hồ. Ngược lại, nếu truy vấn \eqref{eq:flip_query} thỏa được tại một ô
nào đó thì mô hình tìm được khác $\sigma^\star$ tại ô ấy, tức là nghiệm thứ hai.
Nói cách khác, nghiệm là duy nhất khi và chỉ khi mọi ô đều được
$\mathit{KB}_{\mathrm{full}}(\mathcal{P}_0)$ suy ra theo nghĩa của
\eqref{eq:entailment}; đây chính là phép phân loại Criminal/Innocent ở Mục~4.4
áp dụng trên tri thức đầy đủ thay vì tri thức công khai. Phép kiểm tra chỉ xét
các biến ô và không quan tâm đến biến phụ trợ của \texttt{ParityConstraint} hay
\texttt{ConnectedConstraint}, vì giả thiết được truyền vào bộ giải luôn là một
phán quyết trên ô.

Nhờ vậy, thay vì phải liệt kê toàn bộ $2^{|\mathcal{C}|}$ phép gán, hàm chỉ cần
$1 + |\mathcal{C}| - |\mathcal{P}_0|$ lời gọi \texttt{has\_solution}. Cả sáu bản
đồ đều có đúng một phán quyết ban đầu, nên số lời gọi bằng đúng số ô, tức $9$,
$16$ và $25$ lần cho lưới $3\times3$, $4\times4$ và $5\times5$. Nhóm đã chạy lại
\texttt{validate\_puzzle(..., require\_clue\_per\_cell=True)} trên cả sáu bản
đồ với cả \texttt{BruteForceChecker} lẫn \texttt{DPLLChecker} và đều thu được kết
quả hợp lệ, trong khi bộ kiểm thử hiện chỉ gọi hàm này với
\texttt{BruteForceChecker} (\texttt{test\_maps}, \texttt{test\_validation},
\texttt{test\_integration}). Ca \texttt{test\_two\_solutions\_rejected} minh họa
trường hợp lỗi: trên bàn $2\times2$ với hai lời khai ``có đúng một tội phạm'' và
``ô 2, 3 không có tội phạm'', cả ô~0 lẫn ô~1 đều có thể là tội phạm, nên hàm báo
\texttt{not unique: cells [0, 1]}.

\paragraph{Phân biệt với kiểm tra giải không cần đoán}
Nghiệm duy nhất là điều kiện cần nhưng chưa đủ để một bản đồ chơi được theo luật
Griductive. Bước kiểm tra cuối cùng của \texttt{validate\_puzzle} là
\texttt{\_check\_solvable\_without\_guessing}, hoạt động theo cách hoàn toàn
khác. Hàm này tạo một \texttt{GameEngine} thật từ bản đồ và cho một ``người chơi
hoàn hảo'' biết trước đáp án chơi qua engine theo từng vòng: trong mỗi vòng, nó
lần lượt nộp phán quyết đúng cho mọi ô chưa chứng minh theo thứ tự định danh, và
mỗi phán quyết được chấp nhận sẽ lật lời khai mới theo đúng giao thức
\eqref{eq:reveal}. Nếu một vòng trôi qua mà không có phán quyết nào được chấp
nhận trước khi ván hoàn tất, hàm báo lỗi \texttt{Puzzle requires guessing} kèm
danh sách các ô còn lại. Điểm mấu chốt là ở đây engine quyết định chấp nhận dựa
trên $\mathit{KB}_t$ với các lời khai \emph{đã lật}, chứ không phải toàn bộ
$\mathcal{K}$, nên bước này kiểm tra được thứ tự mở thông tin mà bước kiểm tra
tính duy nhất bỏ qua. Bảng~\ref{tab:unique_vs_guess} tóm tắt sự khác biệt giữa
hai bước.

\begin{table}[htbp]
  \centering\small
  \caption{So sánh hai bước kiểm chứng ngữ nghĩa trong \texttt{validate\_puzzle}}
  \label{tab:unique_vs_guess}
  \begin{tabularx}{\textwidth}{l>{\raggedright\arraybackslash}X>{\raggedright\arraybackslash}X}
    \toprule
    & \textbf{\texttt{\_check\_unique\_\allowbreak solution}}
    & \textbf{\texttt{\_check\_solvable\_\allowbreak without\_guessing}} \\
    \midrule
    Tri thức dùng       & Toàn bộ lời khai $\mathcal{K}$ và $\mathcal{P}_0$
                        & Chỉ lời khai đã lật $\mathcal{R}_t$ và $\mathcal{P}_t$ tại từng thời điểm \\
    Giao thức lật mở    & Không xét & Áp dụng đúng như khi chơi, qua \texttt{GameEngine} \\
    Cách gọi bộ giải    & Gọi trực tiếp \texttt{has\_solution}
                        & Gián tiếp qua \texttt{submit\_verdict}, hai truy vấn mỗi lần nộp \\
    Câu hỏi trả lời     & Có tồn tại nghiệm thứ hai không?
                        & Tại mọi thời điểm chưa xong, có ô nào suy ra được từ thông tin công khai không? \\
    Thông báo lỗi       & ``does not satisfy all clues''; ``not unique: cells [...]''
                        & ``requires guessing'' kèm danh sách ô; ``contradict the solution'' \\
    \bottomrule
  \end{tabularx}
\end{table}

Giữa hai bước có một quan hệ một chiều. Nếu bản đồ giải được không cần đoán thì
mọi ô đều đã được suy ra từ một $\mathit{KB}_t$, mà $\mathit{KB}_t$ chỉ là một
phần của $\mathit{KB}_{\mathrm{full}}(\mathcal{P}_0)$ cộng thêm các phán quyết
vốn cũng được suy ra theo cách đó, nên bằng quy nạp mọi ô đều được
$\mathit{KB}_{\mathrm{full}}(\mathcal{P}_0)$ suy ra và nghiệm là duy nhất. Chiều
ngược lại thì không đúng, và ca kiểm thử \texttt{test\_requires\_guessing} là
một phản ví dụ cụ thể. Bàn $2\times2$ ở đây có ba lời khai: hai lời khai công
khai ``có đúng một tội phạm'' và ``ô 2, 3 không có tội phạm'', cùng lời khai
``ô 1 vô tội'' do chính ô~1 nắm giữ. Xét trên toàn bộ lời khai, nghiệm là duy
nhất (ô~0 là tội phạm), nên bước 2 không báo lỗi. Tuy nhiên khi chơi, sau khi ô~2
và ô~3 được chứng minh vô tội thì ô~0 và ô~1 vẫn hòa nhau, trong khi lời khai
phá thế hòa lại chỉ được lật khi ô~1 đã được chứng minh, điều không bao giờ xảy
ra được. Hàm \texttt{\_check\_solvable\_without\_guessing} vì vậy báo
\texttt{requires guessing} với các ô \texttt{[0, 1]}.

\paragraph{Vì sao vẫn giữ cả hai bước}
Dù bước không đoán đã kéo theo tính duy nhất, bước kiểm tra tính duy nhất vẫn
có ích theo hai cách. Thứ nhất, nó chẩn đoán lỗi rõ hơn: một bản đồ có nhiều
nghiệm là lỗi ở chính tập lời khai, còn một bản đồ nghiệm duy nhất nhưng phải
đoán là lỗi ở cách phân bố lời khai giữa các ô, và việc tách hai bước giúp người
thiết kế biết cần sửa ở đâu. Thứ hai, bước 1 là nơi duy nhất bảo đảm đáp án khai
báo thỏa \emph{mọi} lời khai. Vòng lặp của bước không đoán kết thúc ngay khi
\texttt{is\_complete} đúng, nên lời khai được lật bởi phán quyết cuối cùng không
bao giờ được đem ra kiểm tra trong vòng lặp đó; bước 1 bù lại điều này. Cũng vì
bước 1 đã chạy trước, nhánh báo lỗi \texttt{CONTRADICTED} trong bước không đoán
thực tế không thể xảy ra: khi $\sigma^\star$ thỏa mọi lời khai thì nó cũng thỏa
mọi tập con đã lật, nên trạng thái đúng của một ô luôn còn khả dĩ.

Cuối cùng, cần lưu ý rằng một ``vòng'' của bước không đoán không trùng với một
``sóng suy luận'' ở Mục~\ref{sec:reveal_protocol}. Trong cùng một vòng, ảnh chụp
công khai được cập nhật ngay sau mỗi phán quyết được chấp nhận, nên lời khai vừa
lật có thể được dùng cho các ô có định danh lớn hơn ngay trong vòng đó, còn một
sóng chỉ dùng thông tin có ở đầu sóng. Số vòng do hàm này thực hiện vì vậy có
thể nhỏ hơn số sóng ghi trong Bảng~\ref{tab:map_overview}, và hàm chỉ được dùng
để trả lời câu hỏi có giải được hay không chứ không dùng để đo độ khó.
\end{sloppypar}
